\section{Randomized consequences from the Lean-scoped fixed-width theorem}
\label{sec:lean-connectivity}

This section assumes \HLean, the fixed-width exact-counting assertion
specified in Section~\ref{sec:interface}, using only $k=3$.
Put
\[
 f_2(n)=\frac{n(\log_2\log_2(n+16))^2}{\sqrt{\log_2(n+16)}}.
\]
Graphs are simple and loopless, with $m\le cn$ for fixed $c$.
Undirected cycles are counted as edge sets and have at least three
vertices. Opposite arcs are allowed in a digraph. Boundary cases of
constant size are handled directly. The graph and all sampled weights
are constants of a constructed formula, rather than its free inputs.

\begin{theorem}[Randomized sparse connectivity, conditional on \HLean]
\label{thm:lean-randomized}
Assume \HLean. Fix $a,c>0$. Hamiltonian-cycle and Hamiltonian-path
decision, with specified or unspecified path endpoints, and witness
recovery on undirected graphs or digraphs with $m\le cn$ admit
randomized algorithms in time
\begin{equation}\label{eq:lean-time}
 \poly_{a,c}(n)2^{O_{a,c}(f_2(n))}
\end{equation}
and polynomial bit space. The same bounds hold for minimum-cost
Hamiltonian TSP, undirected Steiner tree, feedback vertex set,
connected vertex cover, and connected dominating set, with
nonnegative integer edge or vertex costs at most $n^a$, and for
directed Steiner arborescences with a specified root and direction.

A decision trial has no false positives and, whenever a feasible
object exists, returns a valid witness with probability at least $1/2$.
An optimization trial returns an optimal value and witness with
probability at least $1/2$ whenever the feasible family is nonempty.
Every returned witness is feasible; a failed optimization trial can
return a valid but suboptimal solution or report no solution found.
With $R$ independent trials and the best returned witness, the success
probability is at least $1-2^{-R}$; time is multiplied by $R$ and space
is polynomial in $n+\log(R+1)$. If the manuscript's fixed-linear-density
case of \HBB\ is assumed instead, $f_2(n)$ in~\eqref{eq:lean-time}
can be replaced by $g(n)$.
\end{theorem}

Two trials give success at least $3/4$. A claimed cost bound can be
checked on a returned witness; optimality is a probabilistic guarantee,
not a certificate supplied by every successful feasibility check.
For Hamiltonian TSP the same proof permits signed integer costs of
absolute value at most $n^a$, by the normalization below. The other
optimization problems use the nonnegative-cost conventions just stated.

\subsection{Linear circuits for fixed polynomially bounded weights}
\label{sec:lean-arithmetic}

First consider the population count of $d$ Boolean bits. Place the
bits in the unit column. Replace three live bits in one column by the
sum bit in that column and the carry bit in the next, using a full
adder with two XORs, two ANDs, and one OR. Each replacement reduces
% Sum and carry are explicit functional outputs.
the number of live bits by one. The sum of their column weights is
invariant and equals $d$, so no live wire can lie in a column of
weight exceeding $d$. After at most $d$ replacements there are at
most two bits per column. A final ripple addition has
$O(\log(d+2))$ gates. This gives $O(d+1)$ fan-in-two gates, with
unique auxiliary values. The cases $d=0,1$ are immediate.

Let $S\ge2$, $p\le S$, and let the known constants satisfy
$0\le a_i\le S^a$, with $a$ fixed. We construct a circuit for the
complete integer sum $\sum_{i=1}^p a_ix_i$. Choose
\[
 d=\max\{1,\lfloor\log_2 S/4\rfloor\},\qquad D=2^d,
 \qquad a_i=\sum_{j=0}^{h-1}a_{ij}D^j.
\]
Here $0\le a_{ij}<D$ and $h=O_a(1)$; the constructor obtains $h$
from the actual coefficient bit lengths. For each nonempty bucket
$I_{j,u}=\{i:a_{ij}=u\}$ with $1\le u<D$, compute
$C_{j,u}=\sum_{i\in I_{j,u}}x_i$ by a population counter.
For each $j$ the buckets are disjoint, so these counters together
cost $O(p+D)$ gates, not $O(pD)$. Different digit positions continue
to read the same original bits $x_i$.

Writing $q=\lceil\log_2(p+1)\rceil$, a shift-and-add multiplication
by the constant $u<D$ costs $O(d(q+d+1))=O(\log^2 S)$ gates per
bucket. There are at most $h\min(p,D-1)$ nonempty buckets. Shift
each resulting number by $dj$ and add them using
$Q=\lceil\log_2(p\max_i a_i+1)\rceil+1=O_a(\log S)$ bits.
Partial sums never exceed $p\max_i a_i$, so there is no overflow.
The total gate count is
\begin{equation}\label{eq:lean-weight-size}
 O_a\bigl(p+D\log^2 S+D\log S\bigr)=O_a(S).
\end{equation}
All-zero coefficients and $p=0$ are handled directly. Reading the
$O_a(S\log S)$ coefficient bits takes polynomial time; those bits
do not become free formula variables. Constantly many sums,
comparisons, and cardinality constraints preserve the bound.

An equivalent implementation computes digit subtotals
$T_j=\sum_u uC_{j,u}\le p(D-1)$ and uses carries
\[
 K_0=0,\quad U_j=T_j+K_j,\quad
 e_j=U_j\bmod D,\quad K_{j+1}=\lfloor U_j/D\rfloor.
\]
Inductively $0\le K_j\le p$, hence $U_j\le pD$; each digit step
uses $O(\log S)$-bit arithmetic. The identity
$T_j=e_j+DK_{j+1}-K_j$ telescopes to
\[
 \sum_{j<h}T_jD^j=\sum_{j<h}e_jD^j+K_hD^h.
\]
Keeping the final carry gives the full integer sum. This is the
carry implementation exercised by the ancillary programs; both
implementations satisfy~\eqref{eq:lean-weight-size}.

Signed coefficients are normalized using literals. Set
$C=\sum_{a_i<0}a_i$, and take $\ell_i=x_i$ for $a_i\ge0$ and
$\ell_i=1-x_i$ otherwise. Then
$\sum_i a_ix_i=C+\sum_i|a_i|\ell_i$. Translate the target by $C$.
A complemented literal creates no independent free bit. This also
handles deletion costs $\sum_v c_v(1-u_v)$.

Every gate is encoded by its complete Tseitin equivalence. For example,
$z=x\wedge y$ uses
$(\neg z\vee x)$, $(\neg z\vee y)$, and
$(z\vee\neg x\vee\neg y)$. For $z=x\oplus y$ the four clauses are
\[
\begin{aligned}
 &(x\vee y\vee\neg z),&& (x\vee\neg y\vee z),\\
 &(\neg x\vee y\vee z),&& (\neg x\vee\neg y\vee\neg z).
\end{aligned}
\]
OR follows by De Morgan's law using complemented literals; NOT uses
$(z\vee x)$ and $(\neg z\vee\neg x)$, and constants use units.
There is one variable per gate output. Topological induction gives
exactly one assignment to these variables for each assignment to the
primary bits. Unused declared inputs retain both values. Unit and
binary clauses are permitted by the fixed-width interface, so no
padding variables are needed.

\subsection{Weight layers, isolation, and recovery}
\label{sec:lean-layers}

Fix a root $r$. Use edge bits $x_e$ and cut-color bits $c_v$, and
fix $c_r=0$. For $e=uv$ impose the guarded equalities
\begin{equation}\label{eq:lean-cut-clauses}
 (\neg x_e\vee\neg c_u\vee c_v),\qquad
 (\neg x_e\vee c_u\vee\neg c_v).
\end{equation}
Thus an unselected edge imposes no cut constraint. Population
counters require degree two at every vertex. The circuit above
requires $\sum_e w_ex_e=t$, for fixed polynomially bounded $w_e$.
For $S=n+m+2$, the resulting 3-CNF $\Phi_{G,w,t}$ has
$N,M=O_a(S)$ variables and clauses, and
\begin{equation}\label{eq:lean-layer-count}
 \#\Phi_{G,w,t}
 =\sum_{\substack{F\subseteq E:\deg_F(v)=2\ (v\in V)\\w(F)=t}}
       2^{\cc(V,F)-1}.
\end{equation}
Indeed all arithmetic states are unique, each selected component
has a common color, the root component has color zero, and every
other component has two choices. Degree-two candidates are disjoint
unions of cycles. Modulo two,~\eqref{eq:lean-layer-count} therefore
counts exactly the Hamiltonian cycles in the layer. Its positive
integer value alone does not certify Hamiltonicity.

The size bound is total and amortized: a vertex of degree $d_v$
uses $O(d_v+1)$ gates, and $\sum_vd_v=2m$. No constant bound on
the maximum degree is imposed. More explicitly, if there are $G$
gate outputs, then $N=m+n+1+G$ and at most $4G$ gate clauses,
$2m$ cut clauses, $O(n+m)$ degree-output pins, and $O_a(\log S)$
weight-output pins. Thus $M=O_a(S)$ as well. Pinning any subset of
object bits adds at most $O(S)$ units. In these circuits $G$ bounds
the number of output pins, and the remaining primary bits are
declared, so $M\le D_0N$ for a fixed constant $D_0$.

For completeness, the needed isolation bound follows by conditioning
on all weights except one. For a nonempty family
$\mathcal F\subseteq2^{[p]}$, let every $w_i$ be independent and
uniform in $[W]$. For each $i$, the minimum weight among members
containing $i$ has the form $w_i+A$, while the minimum among members
not containing $i$ is $B$, with $A,B$ independent of $w_i$.
If both types occur, equality is possible at at most one value of
$w_i$. If the global minimum is attained by distinct members, some
coordinate belongs to one and not the other, giving such an equality.
The union bound gives failure probability at most $p/W$.
This is the standard isolation step in Cut\&Count~\cite{CutCount11}.

Reject degree-deficient HC inputs before sampling. For the remaining
inputs take $p=m$ and $W$ the least power of two at least $2m$,
so $2m\le W<4m$ on nontrivial inputs. Sampling takes exactly
$m\log_2 W$ random bits. Scan $0\le t\le nW$ and count each
formula exactly using \HLean, then reduce modulo two. No-HC inputs
always reject. When the HC family is nonempty, it has a uniquely
minimum-weight member with probability at least $1/2$; its layer
is odd and is found by the scan.

Keep the weights and an odd target layer fixed. If the current count
is odd, the counts after setting a chosen object bit to zero and to
one sum to it. Exactly one child count is odd. Query the zero child
and retain the odd branch. After at most $m$ queries the edge set is
fixed and its odd cut multiplicity forces a genuine HC. Cut colors
and all gate outputs remain summed over throughout. Recovery is
deterministic given the odd layer and loses no success probability.
One can check the final cycle directly.

Each formula has a polynomial-length canonical encoding. Even the
unary natural-number encoding used by the fixed-width interface
has length $O(NM+N)=\poly(n)$ here; a binary clause-list length is
not being assumed for that interface. Each exact answer has at most
$N+1$ bits. There are only polynomially many layers and recovery
queries. Sequential calls to the single fixed $k=3$ machine therefore
give~\eqref{eq:lean-time} and polynomial bit space. Under the
manuscript's stronger fixed-density statement, use $M\le D_0N$
and obtain exponent $g(n)$. These are distinct assumptions.

\subsection{Optimization and other connectivity predicates}
\label{sec:lean-objects}

More generally, let $x\in\{0,1\}^p$, $p=O(n)$, encode an object,
and let a linear-size 3-CNF $\Psi(x,z)$ satisfy
\begin{equation}\label{eq:lean-parity-interface}
 \#\{z:\Psi(x,z)\}\equiv[x\text{ is a true object}]\pmod2.
\end{equation}
Isolate the object bits $x$, not their auxiliary cut assignments.
Adding a fixed-weight sum and target preserves linear size. The same
layer scan and odd-branch recovery apply. If $p=0$, query the one
empty object directly.

For optimization with a linear integer cost
$C(x)=C_0+\sum_i d_ix_i$, whose coefficients and range are fixed
polynomially bounded quantities, take $P=pW+1$ and use the statistic
\begin{equation}\label{eq:lean-lexicographic}
 A(x)=P C(x)+\sum_{i=1}^p w_ix_i.
\end{equation}
The isolation-weight difference between two objects is at most $pW$,
so $A$ orders costs first. Its coefficients and full range remain
polynomially bounded. Scan its range in increasing order, including
the known constant shift if signed literals are used, and recover
from the first odd layer. Apply isolation to the nonempty family of
originally optimal objects. With probability at least $1/2$ its
minimum is unique, and that object is the unique global minimum of
$A$. Thus the first odd layer then yields an original optimum.
Off this event, all returned objects remain feasible but can be
suboptimal. Independent repetition and retaining the smallest
original cost proves the stated probability bound. This avoids
assuming that a sequence of unamplified adaptive threshold decisions
has the same success probability as a single query.

Here are the predicates realizing~\eqref{eq:lean-parity-interface}.
All cut equalities are guarded as in~\eqref{eq:lean-cut-clauses}.
Arithmetic and size constraints use Section~\ref{sec:lean-arithmetic}.
\begin{enumerate}
\item For an undirected Hamiltonian $s$--$t$ path, the endpoint
degrees are one and the remaining degrees two; fix the root color
at $s$. The two odd-degree vertices lie in the same component.
Every candidate is an $s$--$t$ path plus disjoint cycles, so the
root-fixed cut parity retains precisely the spanning paths.
Enumerate unordered endpoint pairs if unspecified; a one-vertex
path is treated directly.

\item In a digraph use selected indegree and outdegree one at every
vertex for HC. For a path use $(0,1)$ at its start, $(1,0)$ at
its end, and $(1,1)$ elsewhere. Cut constraints use the underlying
undirected incidences. Weak connectivity makes the candidate a single
directed cycle or path. Enumerate ordered endpoint pairs when needed.

\item For an undirected Steiner tree isolate $(U,F)$, require
$T\subseteq U$, both endpoints of every selected edge in $U$,
and $|F|=|U|-1$. Fix the color of a root in the nonempty terminal
set $T$, and fix colors outside $U$ to zero. The cut multiplicity is
$2^{\cc(U,F)-1}$; oddness gives connectedness and then the cardinality
equality gives a tree. If $T=\varnothing$, handle the empty-tree
convention directly. With nonnegative costs, minimum connected
Steiner subgraphs have an optimum tree.

\item For connected vertex cover isolate $U$, require
$u_v\vee u_w$ on every input edge, and let induced-edge bits be
functional ANDs $u_vu_w$. For connected domination instead require
the OR of the selected bits in every closed neighborhood to be true.
Such OR chains have $O(n+m)$ inputs in total. For either problem,
fix a selected root, set colors outside $U$ to zero, and impose
same-color constraints on induced edges. Oddness is exactly
connectedness. Enumerate selected roots; an empty allowed solution
is checked separately. Choose a root in one optimal solution for
the success-probability analysis; no simultaneous success for all
roots is required.

\item For FVS isolate $(U,M_0)$, where $U$ is retained and
$M_0\subseteq U$ is a set of markers. Let $F=E[U]$ be functional
AND outputs and impose $|F|+|M_0|=|U|$. Fix colors of markers and
of unselected vertices to zero, with no additional unconditional
root pin. The multiplicity is $2^z$, where $z$ counts selected
components without a marker. Oddness gives $|M_0|\ge\cc(U,F)$.
The elementary bound $|F|\ge|U|-\cc(U,F)$ and the equality give
the reverse inequality. Thus each component has one marker and
is a tree. Conversely such a marked forest has exactly one cut.
The deleted set $V\setminus U$ is an FVS. Charge its cost only to
deletions but isolate both $U$ and $M_0$; omitting marker bits from
isolation could cancel all representations of a valid forest.
Every retained forest admits a marker in each component, without
changing the original cost, so optimization over these representations
has the same optimum as the deletion problem.
Empty $U=M_0=\varnothing$ has one auxiliary assignment.

\item For an inward directed Steiner arborescence, isolate $(U,F)$
and explicitly impose
\[
 u_r=1,\qquad u_v=1\quad(v\in T).
\]
These terminal requirements are unit clauses. The root has selected
outdegree zero, other selected vertices outdegree one, and unselected
vertices outdegree zero; selected arcs have both endpoints selected.
Fix $c_r=0$ and all colors outside $U$ to zero, and use weak cut
consistency on selected arcs. The arc count is $|U|-1$.
Oddness gives weak connectivity, so the underlying multigraph is a
tree. Following unique outgoing arcs cannot cycle and must reach
the root. Conversely an inward tree has one root-fixed cut.
Reverse all arcs for the outward version. Any feasible reachability
subgraph has such an arborescence, and for nonnegative costs removing
extra arcs does not increase the cost.
\end{enumerate}

Each construction uses $O(n+m)$ primary bits and gate inputs in total;
the degree counters, functional induced edges, neighborhood ORs, and
fixed number of statistics all have linear size. These observations,
isolation, and sequential fixed-width counting prove
Theorem~\ref{thm:lean-randomized}.

\subsection{Deterministic sublinear-bit modular counting}
\label{sec:lean-modular}

\begin{theorem}[HC residues, conditional on \HLean]
\label{thm:lean-modular}
Assume \HLean. On simple undirected graphs with $m\le cn$, given
an integer precision $b\ge1$, one can deterministically compute
$\#\mathrm{HC}(G)\bmod2^b$. For every supplied precision sequence
$b=b(n)=o(n)$, writing $t=b+v_2(n)$, the time is
\begin{equation}\label{eq:lean-modular-time}
 \poly_c(n+b)\,
 2^{O_c(f_2(n)+t\log_2(en/t))}=2^{o(n)},
\end{equation}
with polynomial bit space, for all sufficiently large $n$.
The fixed-density case of \HBB\ permits $f_2$ to be replaced by $g$.
The precision is an explicit input; any cost of generating it must be
included if it is generated by an algorithm instead.
\end{theorem}

Here $v_2(n)$ is the exponent of two in $n$. We give a uniform
algorithm for arbitrary $b$, with a finite-sum running-time bound
below, and derive~\eqref{eq:lean-modular-time} when $b=o(n)$.
Handle $n<3$ directly and fix a root $r$ otherwise.

Use edge bits $x_e$ with degree-two constraints. For each potential
matrix position in
\[
 \mathcal P=\{(v,v):v\in V\}\cup
     \{(u,v),(v,u):uv\in E\},\qquad q=|\mathcal P|=n+2m,
\]
introduce a bit $z_{ij}$, and require exactly one selected position
in each row and column. A selected off-diagonal position requires
its edge bit $x_e=1$. The degree counters and row/column counters
have total incidence $O(n+m)$.

For each nonroot diagonal add a multiplicity bit implying $z_{vv}$.
It has one assignment when $z_{vv}=0$ and two when $z_{vv}=1$.
For the root add two bits, each implying $z_{rr}$, and forbid their
simultaneous truth. They give one or three assignments. The resulting
unsigned 3-CNF $\Psi$ has $N,M=O(n+m)$, with all ordinary gate states
unique. It realizes the absolute determinant terms of
\[
 A_F=L_F+e_re_r^{\mathsf T}:
 \quad (A_F)_{vv}=2+[v=r],\qquad
 (A_F)_{uv}=-x_{uv}\quad(u\ne v)
\]
for every selected 2-factor $F$, with $x_{uv}=0$ on nonedges.
The matrix-tree identity gives
$\det A_F=0$ if $F$ is disconnected and $\det A_F=n$ if it is a
Hamiltonian cycle. Therefore the signed sum of assignments of
$\Psi$ is $n\,\#\mathrm{HC}(G)$.

To evaluate that sum without a large sign circuit, define on every
Boolean position vector
\[
 Q(z)=\sum_{\substack{(i,j),(k,l)\in\mathcal P\\i<k,\ j>l}}
             z_{ij}z_{kl},\qquad
 f(z)=(-1)^{Q(z)+\sum_{i\ne j}z_{ij}}.
\]
On a selected permutation this is exactly its sign times the signs
of its negative matrix entries. For $S\subseteq\mathcal P$ put
\begin{equation}\label{eq:lean-mobius}
 c(S)=\sum_{R\subseteq S}(-1)^{|S|-|R|}f(1_R).
\end{equation}
Inclusion--exclusion gives $f(1_X)=\sum_{S\subseteq X}c(S)$:
after exchanging the sums, the coefficient for $R\subseteq X$ is
$(1-1)^{|X\setminus R|}$, leaving only $R=X$.

In the Boolean integer polynomial ring, $f$ is a product of factors
$1-2z_iz_j$ for inversions and $1-2z_i$ for off-diagonal positions.
A term whose support has size $s$ uses at least $\lceil s/2\rceil$
nonconstant factors. Reducing $z_i^2=z_i$ and adding terms with the
same support preserves their divisibility. Uniqueness of the integer
multilinear expansion, or~\eqref{eq:lean-mobius}, thus gives
\[
 2^{\lceil|S|/2\rceil}\mid c(S).
\]
Modulo $2^t$, every support with $|S|\ge2t-1$ vanishes. Let
$\Psi_S$ be $\Psi$ with $z_{ij}=1$ pinned for positions in $S$.
Exchanging the finite sums now yields
\begin{equation}\label{eq:lean-modular-expansion}
 n\,\#\mathrm{HC}(G)\equiv
 \sum_{|S|\le\min(q,2t-2)}c(S)\#\Psi_S\pmod{2^t}.
\end{equation}
Each query is still a linear-size 3-CNF. Conflicting row or column
pins can be skipped since they have count zero.

Write $h=v_2(n)$ and $n=2^hu$ with $u$ odd. Compute the least
nonnegative residue $R_0$ in~\eqref{eq:lean-modular-expansion} at
$t=b+h$. It is divisible by $2^h$, and the answer is
\[
 (R_0/2^h)u^{-1}\pmod{2^b}.
\]
This divides in the integers before inverting an odd number;
it never inverts an even $n$ modulo a power of two.

Computing each coefficient by its $2^{|S|}$-term definition gives
the uniform bound
\[
 \poly_c(n+b)\,2^{O_c(f_2(n))}
 \sum_{s=0}^{\min(q,2t-2)}\binom qs2^s.
\]
If $t=o(n)$ and $q=\Theta_c(n)$, the sum is
$2^{O_c(t\log_2(en/t))}$. Since $v_2(n)=O(\log n)$, supplied
$b=o(n)$ makes $t=o(n)$, and
$(t/n)\log(en/t)\to0$, proving the stated subexponential bound.
Enumerate $S$ and $R$ sequentially, storing one formula, one exact
count, and a residue accumulator. This uses polynomial space.
The argument invokes only explicit fixed-width formulas, not a
tensor-contraction primitive.

\subsection{Comparison and scope of the two inputs}

Cut\&Count, isolation, and functional Tseitin encoding are established
frameworks~\cite{CutCount11}. The randomized algorithms here are their
composition with fixed-width exact counting, after the explicit linear
encoding just proved. Its role here is a proved size and
multiplicity bound.

The main bound uses the narrower \HLean\ exponent $f_2$, not the
manuscript's sharper fixed-density bound. The reductions themselves
are not claimed to have been verified in Lean. Polynomially bounded
costs mean $O(\log n)$ bits with a fixed exponent on their numerical
size; the $n^{9/10}$-bit costs of the deterministic tensor results are
outside this statement. Randomized feasibility and probabilistic
optimality here do not replace the deterministic results of
Sections~\ref{sec:determinants}--\ref{sec:directed-routes}.
Similarly, a positive HC count may be divisible by $2^b$, so a zero
residue does not imply infeasibility. Sublinear-bit residues are not
full integer counting on arbitrary sparse graphs. General
deterministic HC decision, full exact counting, and the long-cost
range remain \HTC-dependent results in this paper; no impossibility
of deriving them from a weaker hypothesis is asserted.
